% ============================================================
%  Emergent Gauge Structure from a 3D Brane in 4D Embedding
% ============================================================
\documentclass[11pt,a4paper]{article}

% ---------- Packages ----------
\usepackage[T1]{fontenc}
\usepackage[utf8]{inputenc}
\usepackage{lmodern}
\usepackage{microtype}

\usepackage{amsmath,amssymb,amsfonts,mathtools}
\usepackage{bm}
\usepackage{physics} % bras/kets, derivatives (optional)

\usepackage{graphicx}
\usepackage{subcaption}

\usepackage{hyperref}
\hypersetup{
  colorlinks=true,
  linkcolor=blue,
  citecolor=blue,
  urlcolor=blue
}

\usepackage{geometry}
\geometry{margin=2.5cm}

% ---------- Macros ----------
\newcommand{\R}{\mathbb{R}}
\newcommand{\C}{\mathbb{C}}
\newcommand{\T}{\mathbb{T}}
\providecommand{\dd}{\mathrm{d}}
\providecommand{\ii}{\mathrm{i}}
\providecommand{\e}{\mathrm{e}}

% Indices:
%   i,j,k,... intrinsic spatial indices (1..d)
%   A,B,... embedding indices (1..N)
%   mu,nu,... spacetime indices (0..3)

\newcommand{\X}{\mathbf{X}}                 % embedding map
\newcommand{\x}{\mathbf{x}}                 % intrinsic coordinate (space)
\newcommand{\uvec}{\mathbf{u}}              % displacement field
\newcommand{\vvec}{\mathbf{v}}              % velocity field
\newcommand{\A}{\mathbf{A}}                 % gauge potential (spatial)
\newcommand{\F}{\mathbf{F}}                 % field strength / curvature
\newcommand{\Lag}{\mathcal{L}}              % Lagrangian density
\newcommand{\Ham}{\mathcal{H}}              % Hamiltonian density

% Berry / Wilczek–Zee
\newcommand{\bbA}{\mathcal{A}}              % non-Abelian connection
\newcommand{\bbF}{\mathcal{F}}              % non-Abelian curvature
\newcommand{\ab}{a}                         % U(1) Berry connection scalar components a_mu
\newcommand{\covD}{D}                       % covariant derivative symbol

% QED-like notation
\newcommand{\psibar}{\bar{\psi}}
\newcommand{\gamm}{\gamma}

% ---------- Title ----------
\title{Emergent Gauge Structure and Particle-Like Modes\\
from a 3D Brane Embedded in 4D Space}
\author{Author: \textit{(your name)}}
\date{\today}

\begin{document}
\maketitle

\begin{abstract}
We propose a deterministic wave model in which the fundamental substrate is a three-dimensional elastic brane
embedded in $\mathbb{R}^4$, evolved in an external time parameter.
Electromagnetism is treated as an \emph{emergent} gauge structure that appears when the real brane dynamics is
projected onto a narrowband, locally transported polarization (or near-degenerate internal) subspace, yielding
a U(1) Berry connection and, for higher-rank subspaces, a Wilczek--Zee connection.
Electron-like excitations are treated as localized, self-consistent nonlinear modes of the same substrate.
We present (i) the continuous hyperelastic formulation and its discrete spring-lattice approximation,
(ii) the narrowband multiscale ansatz and the induced connections, and
(iii) simulation protocols that compute gauge-invariant curvatures and test Maxwell-type identities and
confinement stability in a reference implementation.
\end{abstract}

\tableofcontents
\newpage

% ============================================================
\section{Scope and Claims}
\label{sec:scope}

This paper develops a \emph{mechanistic} link between a concrete brane substrate model and effective gauge
structures obtained by mode-space projection. The scope is intentionally narrow: we focus on (a) the emergence
of a U(1) (and optionally Wilczek--Zee) connection from narrowband mode-frame transport, and (b) the existence
and numerical stability of localized ``electron-like'' excitations in the same substrate.

\subsection*{Core claims}
\begin{enumerate}
  \item \textbf{Substrate model (defined precisely).} The fundamental degrees of freedom are an embedding map
        $\X(\x,t):\Omega\subset\mathbb{R}^3\to\mathbb{R}^4$ with a strictly isotropic hyperelastic energy
        $W(g;g^0)$ built from the induced metric $g_{ij}=\partial_i\X\cdot\partial_j\X$ relative to an isotropic
        reference metric $g^0_{ij}=\ell_0^2\delta_{ij}$. A single dimension-free parameter
        $\alpha=\ell_0/h_\star$ controls uniform pre-stress via boundary-enforced geometric scale $h_\star$.
  \item \textbf{Connections are kinematical (and require an ordered narrowband sector).} For narrowband solutions,
        projecting the real brane fields onto a locally transported $n$-dimensional mode subspace produces a covariant derivative
        $\covD_\mu=\partial_\mu+\ii a_\mu\mathbf{1}_n+\ii\bbA_\mu$, where $a_\mu$ is the U(1) Berry connection and
        $\bbA_\mu$ is a Wilczek--Zee connection when $n>1$. These connections encode \emph{basis transport} of the retained subspace
        and are not assumed as additional fundamental fields. Crucially, the reduction is \emph{not} the ``phase of arbitrary vibrations'':
        it presupposes a coherent, narrowband carrier (a macroscopic complex envelope) whose phase is stiff on the envelope scale.
        A broadband ``vibration soup'' generically destroys adiabaticity and produces noisy frame transport rather than a clean gauge field.
  \item \textbf{Localized modes are posed as a self-adjoint eigenproblem plus nonlinearity.} Linearization about a
        stationary, spherically symmetric background yields a self-adjoint operator under the mass-weighted
        inner product. Standard spectral results then provide orthonormality/completeness and justify spherical
        harmonic angular structure as a consequence of $SO(3)$ symmetry (not as an imposed constraint).
  \item \textbf{Testable numerical consequences.} The paper specifies gauge-invariant lattice estimators of
        curvature (plaquette phases) and concrete validation tests (Maxwell identities, polarization counting,
        confinement stability and parameter dependence) for a discrete reference implementation.
\end{enumerate}

\subsection*{Non-claims (boundaries of the present work)}
\begin{itemize}
  \item We do \emph{not} claim a completed derivation of the Standard Model, quantum statistics, or measurement
        theory. The Berry/Wilczek--Zee machinery is used here only as a geometrical reduction tool for a
        narrowband sector of the substrate dynamics.
  \item We do \emph{not} assume a point charge or an independent electromagnetic field at the fundamental level.
        Any effective Maxwell term must arise from coarse-graining; Section~\ref{sec:maxwell_term} states what
        must be derived and what is currently a program.
  \item Relativistic invariance is not postulated at the substrate level (time is an external parameter); only
        effective, sector-specific wave propagation properties are compared to relativistic phenomenology.
\end{itemize}


% ============================================================
\section{Conceptual Model: Continuous 3D Brane in 4D Embedding}

\label{sec:conceptual_model_continuous}
% ============================================================

\subsection{Ontology and kinematics (time as parameter)}
Let $\Omega\subset\mathbb{R}^3$ be the intrinsic (material) brane domain.
We use either a periodic cell $\Omega=\mathbb{T}^3$ or a bounded cube with
clamped boundary points. Time $t\in\mathbb{R}$ is an external evolution
parameter, not an additional coordinate of the brane.

The fundamental field is the embedding map
\begin{equation}
  \mathbf{X}(\mathbf{x},t)\in\mathbb{R}^4,
  \qquad \mathbf{x}=(x^1,x^2,x^3)\in\Omega,\quad t\in\mathbb{R}.
\end{equation}
We write $\mathbf{X}=(X^1,X^2,X^3,X^4)$ and use the Euclidean dot product
in $\mathbb{R}^4$.
All admissible variations are periodic (for $\mathbb{T}^3$) or satisfy the
appropriate boundary conditions (for clamped boundaries).

% ------------------------------------------------------------
\subsection{Spherical material coordinates (coordinate change only)}
\label{subsec:spherical_material_coords}
% ------------------------------------------------------------

\paragraph{Important: this is only a coordinate change.}
All dynamical statements remain those of the coordinate-free formulation
\(\X:\Omega\times\R\to\R^4\) with induced metric \(g_{ij}=\partial_i\X\cdot\partial_j\X\).
We now introduce spherical coordinates on the \emph{intrinsic/material} domain
\(\Omega\subset\R^3\) as a convenient chart for localized excitations.
This does \emph{not} constrain the embedding \(\X\): the components \(X^1,X^2,X^3,X^4\)
remain general functions of the intrinsic coordinates.

\paragraph{Change of intrinsic coordinates.}
Let \(\x=(x^1,x^2,x^3)\) be Cartesian material coordinates on \(\Omega\).
On a local patch (away from the axis singularities) we use
\[
q^a=(r,\theta,\phi),\qquad
x^1=r\sin\theta\cos\phi,\;\;
x^2=r\sin\theta\sin\phi,\;\;
x^3=r\cos\theta,
\]
with \(r\ge 0\), \(\theta\in[0,\pi]\), \(\phi\in[0,2\pi)\).
The Jacobian is
\[
\dd^3x = J(q)\,\dd r\,\dd\theta\,\dd\phi,
\qquad
J(q)=r^2\sin\theta.
\]

\paragraph{Euclidean material metric in spherical coordinates.}
The Euclidean metric \(\delta\) on the material space \(\R^3\) has components
\[
\delta_{ab}(q)=\mathrm{diag}\!\big(1,\; r^2,\; r^2\sin^2\theta\big),
\qquad a,b\in\{r,\theta,\phi\}.
\]
Thus statements like \(g^0_{ij}\propto \delta_{ij}\) remain \emph{isotropic} and
coordinate-independent; only the component form changes under the coordinate map.

\paragraph{Embedding map and induced metric in the spherical chart.}
We now write the same embedding field as
\[
\X=\X(q,t)=\big(X^1(q,t),X^2(q,t),X^3(q,t),X^4(q,t)\big)\in\R^4.
\]
The induced metric is still
\[
g_{ab}(q,t)=\partial_a\X(q,t)\cdot\partial_b\X(q,t),
\qquad a,b\in\{r,\theta,\phi\}.
\]
No approximation is made here; this is purely a change of variables.

\paragraph{Optional: static gauge (transverse height field) for later use.}
For later sections focused on transverse confinement via the geometric nonlinearity in \(X^4\),
it is convenient to consider the \emph{static gauge} as a special case of the general embedding:
we identify the first three embedding components with the standard Euclidean chart of the
material space and treat the fourth component as a scalar height field,
\[
(X^1,X^2,X^3)(q,t) = (x^1(q),x^2(q),x^3(q)),
\qquad
X^4(q,t)=u(q,t).
\]
In the spherical material chart this means
\[
(x^1,x^2,x^3)=(r\sin\theta\cos\phi,\; r\sin\theta\sin\phi,\; r\cos\theta),
\qquad
\X(q,t)=\big(x^1(q),x^2(q),x^3(q),u(q,t)\big).
\]
This is \emph{not} imposed in the model; it is merely a useful representation when in-brane
deformations are negligible compared to transverse excitation. In this gauge the induced metric
simplifies to the geometric-nonlinearity form
\[
g_{ab}(q,t)=\delta_{ab}(q) + \partial_a u(q,t)\,\partial_b u(q,t),
\qquad a,b\in\{r,\theta,\phi\},
\]
so all transverse gradients feed back into the stretch energy \(W(g;g^0)\) through \(g_{ab}\).

\paragraph{Reference metric, equilibrium scale, and coupling parameter.}
In Cartesian material coordinates we use the isotropic reference metric
\(g^0_{ij}=\ell_0^2\delta_{ij}\). In the spherical chart this becomes
\[
g^0_{ab}(q)=\ell_0^2\,\delta_{ab}(q)
=\ell_0^2\,\mathrm{diag}\!\big(1,\; r^2,\; r^2\sin^2\theta\big).
\]
If boundary conditions enforce an equilibrium geometric scale \(h_\star\) (uniform in all
directions), the corresponding equilibrium metric reads
\[
g^\star_{ab}(q)=h_\star^2\,\delta_{ab}(q).
\]
We retain the single, direction-free coupling knob
\[
\alpha := \frac{\ell_0}{h_\star}\in(0,1],
\qquad
\varepsilon := \alpha^{-1}-1=\frac{h_\star-\ell_0}{\ell_0}.
\]
This parameter is unchanged by the coordinate choice.

% ------------------------------------------------------------
\subsection{Continuous Lagrangian in spherical material coordinates}
\label{subsec:continuous_lagrangian_spherical}
% ------------------------------------------------------------

\paragraph{Coordinate-free Lagrangian density (unchanged physics).}
Let \(\rho_m\) be the brane mass density. The Lagrangian density is
\[
\Lag(\X,\partial_t\X,\nabla\X)=\frac{\rho_m}{2}\,|\partial_t\X|^2 - W(g;g^0),
\qquad
g_{ab}=\partial_a\X\cdot\partial_b\X,
\]
where \(W\) is any isotropic hyperelastic energy density minimized at \(g=g^0\)
(equivalently: \(W\) depends only on invariants / principal stretches of
\(\widehat g=(g^0)^{-1}g\)). This matches the rest-length spring family in the
small-strain limit; the coordinate choice does not modify \(W\).

\paragraph{Action in spherical material coordinates.}
Using \(\dd^3x = J(q)\,\dd r\,\dd\theta\,\dd\phi\) with \(J=r^2\sin\theta\), the action is
\[
S[\X]=\int \dd t\int_{\Omega}\dd^3x\;\Lag
=\int \dd t\int_{\Omega_q}\dd r\,\dd\theta\,\dd\phi\; J(q)\,
\Big(\frac{\rho_m}{2}|\partial_t\X|^2 - W(g;g^0)\Big),
\]
where \(\Omega_q\) is the intrinsic domain expressed in \((r,\theta,\phi)\).

\paragraph{Euler--Lagrange equations (same dynamics, component form only).}
Define
\[
S^{ab}:=\frac{\partial W}{\partial g_{ab}},
\qquad
P_A^{\;a}:=\frac{\partial W}{\partial(\partial_a X^A)} = 2\,S^{ab}\,\partial_b X^A,
\qquad A=1,\dots,4.
\]
Then the equations of motion in the spherical chart read
\[
\rho_m\,\partial_{tt}X^A
=\frac{1}{J(q)}\,\partial_a\!\Big(J(q)\,P_A^{\;a}\Big),
\qquad A=1,\dots,4,
\qquad J(q)=r^2\sin\theta.
\]
This is exactly the Cartesian divergence form \(\rho_m \partial_{tt}X^A=\partial_i P_A^{\;i}\),
written in spherical material coordinates via the standard Jacobian factor.

\subsection{Isometric geometry: induced metric and reference metric}
The induced metric on the brane is
\begin{equation}
  g_{ij}(\mathbf{X})=\partial_i\mathbf{X}\cdot\partial_j\mathbf{X},
  \qquad i,j=1,2,3,
  \label{eq:induced_metric}
\end{equation}
which is manifestly invariant under rigid motions of the embedding
$\mathbf{X}\mapsto Q\mathbf{X}+a$ with $Q\in O(4)$.

We choose an \emph{isotropic reference metric}
\begin{equation}
  g^0_{ij}=\ell_0^2\,\delta_{ij},
  \label{eq:reference_metric}
\end{equation}
where $\ell_0>0$ is the continuum analog of the discrete rest length.
Define the mixed relative metric
\begin{equation}
  \widehat g := (g^0)^{-1}g,
  \label{eq:relative_metric}
\end{equation}
with eigenvalues $\lambda_a>0$ and principal stretches
$\sigma_a=\sqrt{\lambda_a}$ ($a=1,2,3$).

\paragraph{Coupling / pre-stretch parameter (continuous).}
Let the boundary conditions enforce a ground-state geometric spacing $h_\star$,
i.e. an equilibrium embedding with $g^\star_{ij}=h_\star^2\delta_{ij}$.
We define the single coupling knob
\begin{equation}
  \alpha := \frac{\ell_0}{h_\star}\in(0,1],\qquad
  \varepsilon := \alpha^{-1}-1=\frac{h_\star-\ell_0}{\ell_0}.
  \label{eq:coupling_alpha}
\end{equation}
$\alpha<1$ corresponds to \emph{uniform pre-stretch / pre-tension}:
the brane is held at spacing $h_\star$ while its stress-free metric would be
$\ell_0^2\delta_{ij}$. This pre-stress is the clean control parameter that
sets the strength of transverse--lateral coupling via the geometric nonlinearity
of $g_{ij}(\mathbf{X})$.

\subsection{Constitutive law (Option 2): St.\ Venant--Kirchhoff (Green--Lagrange strain)}
\label{subsec:stvk}

We adopt a continuum \emph{hyperelastic} law that is strictly isotropic and contains \emph{no microstructure}:
the constitutive response depends only on invariants of the induced metric relative to an isotropic reference metric.
This is the continuum analogue of the ``quadratic metric'' approach and is sufficient for deriving a self-adjoint
eigenproblem on a spherically symmetric background.

\paragraph{Relative metric and Green--Lagrange strain.}
Define the relative metric (Cauchy--Green analogue)
\[
C := (g^0)^{-1}g,
\qquad
E := \frac12(C-I).
\]
Here \(E\) is the Green--Lagrange strain tensor.

\paragraph{StVK energy density.}
We use the St.\ Venant--Kirchhoff strain energy density
\begin{equation}
W(g;g^0)=\mu\,\mathrm{tr}(E^2)+\frac{\lambda}{2}\,\big(\mathrm{tr}\,E\big)^2,
\label{eq:stvk_energy}
\end{equation}
with Lam\'e parameters \(\mu>0\) and \(\lambda\) (units: Pa). Using Poisson ratio \(\nu\in(0,1/2)\),
\[
\nu=\frac{\lambda}{2(\lambda+\mu)},
\qquad
\lambda=\frac{2\mu\nu}{1-2\nu}.
\]

\paragraph{Continuum analogue of pre-tension (the \(\alpha\) knob).}
The reference metric \(g^0=\ell_0^2\delta\) encodes the \emph{stress-free} geometric scale.
If boundary conditions enforce a ground-state scale \(h_\star\) (uniform), then
\[
\alpha:=\frac{\ell_0}{h_\star}\in(0,1]
\]
measures the mismatch between the stress-free scale and the held scale. For \(\alpha<1\) the brane is uniformly
\emph{pre-stressed} (continuum ``pre-tension''), which controls how strongly transverse gradients backreact
on in-brane stresses through the geometric nonlinearity of \(g_{ij}(\X)\).

\subsection{Lagrangian, action, and equations of motion}
Let $\rho_m>0$ be the brane mass density (kg/m$^3$). The Lagrangian density is
\begin{equation}
  \mathcal{L}(\mathbf{X},\partial_t\mathbf{X},\nabla\mathbf{X})
  =\frac{\rho_m}{2}\,|\partial_t\mathbf{X}|^2
  -W\!\left(g(\mathbf{X});g^0\right),
  \label{eq:lagrangian_density}
\end{equation}
and the action is
\begin{equation}
  S[\mathbf{X}] = \int dt \int_{\Omega} d^3x\;\mathcal{L}.
  \label{eq:action}
\end{equation}

Define the (second Piola-type) stress and the first Piola-type tensor
\begin{equation}
  S^{ij}:=\frac{\partial W}{\partial g_{ij}},
  \qquad
  P_{Ai}:=\frac{\partial W}{\partial(\partial_i X^A)}
         =2\,S^{ij}\,\partial_j X^A,
  \qquad A=1,\dots,4.
  \label{eq:piola_stress}
\end{equation}
Hamilton’s principle $\delta S=0$ yields the substrate PDE
\begin{equation}
  \rho_m\,\partial_{tt}X^A = \partial_i P_{Ai},
  \qquad A=1,\dots,4,
  \label{eq:continuous_eom}
\end{equation}
with boundary terms vanishing under periodicity or satisfied by the chosen
(clamped) boundary conditions.

\subsection{Linearized isotropic wave propagation and nonlinear coupling}
Linearize about a uniform ground state
$\mathbf{X}^\star(\mathbf{x})=(h_\star x^1,h_\star x^2,h_\star x^3,0)$ and write
$\mathbf{X}=\mathbf{X}^\star+\mathbf{u}$ with small $\mathbf{u}(\mathbf{x},t)\in\mathbb{R}^4$.
Then, to leading order,
\begin{equation}
  \rho_m\,\partial_{tt}\mathbf{u} \;=\; \mathcal{C}(\alpha,\lambda,\mu)\,\Delta\mathbf{u}
  \;+\;\mathcal{O}(\|\nabla\mathbf{u}\|^2),
  \qquad
  c^2=\frac{\mathcal{C}(\alpha,\lambda,\mu)}{\rho_m}.
  \label{eq:linear_wave}
\end{equation}
Here \(\mathcal{C}(\alpha,\lambda,\mu)\) is the effective small-amplitude stiffness about the pre-stressed background.
In practice we calibrate parameters so that the relevant narrowband propagation speed equals the observed \(c\).
This is \emph{rotationally invariant in the intrinsic 3D} because the operator is
the standard Laplacian $\Delta$.

The \emph{dimensional coupling} is entirely geometric: since
$g_{ij}=\partial_i\mathbf{X}\cdot\partial_j\mathbf{X}$ contains cross-terms
between all embedding components, gradients in $X^4$ feed into the in-brane
metric, hence into $W$, hence into the in-brane stresses. The magnitude of the
resulting coupling in the small-amplitude regime is controlled by the pre-stretch
parameter $\alpha$ (or equivalently $\varepsilon$) through the background stress.

\subsection{Electron-like modes as a self-adjoint eigenproblem (no ad-hoc constraints)}
\label{subsec:selfadjoint_eigenproblem}

In textbook derivations of atomic orbitals, constraints such as periodicity, orthonormality, completeness,
and regularity are often presented as if they must be imposed externally. In our brane model we do \emph{not}
modify the governing equations to ``bake in'' such constraints. Instead we ensure two structural properties:

\begin{enumerate}
\item the linearized operator about the chosen background is \emph{self-adjoint} under the correct inner product,
\item the background configuration used for linearization is \emph{spherically symmetric}.
\end{enumerate}

Then orthonormality, completeness, periodicity, and well-behavedness follow automatically as standard results
of spectral theory.

\paragraph{Self-adjoint operator and inner product.}
Let \(\X^\star\) be a stationary background embedding and write \(\X=\X^\star+\eta\).
The second variation of the action defines a linear operator \(\mathcal{L}\) such that
\[
\rho_m\,\partial_{tt}\eta=-\mathcal{L}\eta.
\]
Under periodic boundaries or clamped boundaries, \(\mathcal{L}\) is symmetric with respect to the mass-weighted
inner product
\begin{equation}
\langle \eta,\zeta\rangle :=\int_{\Omega_q}\rho_m\,(\eta\cdot\zeta)\,J(q)\,\dd r\,\dd\theta\,\dd\phi,
\qquad
J(q)=r^2\sin\theta,
\label{eq:inner_product_mass}
\end{equation}
and thus admits a real spectrum with an orthonormal eigenbasis (after specifying the operator domain).

\paragraph{Spherical symmetry and spherical harmonics.}
If \(\X^\star\) is spherically symmetric in intrinsic coordinates, \(\mathcal{L}\) commutes with intrinsic
rotations \(SO(3)\). Consequently eigenmodes can be chosen to transform in irreducible representations of \(SO(3)\),
so the angular dependence is diagonalized by the \(SO(3)\) Casimir operator \(L^2\).
In spherical coordinates this Casimir has the coordinate form of the spherical Laplacian,
This appearance of the Laplace--Beltrami operator is purely a consequence of $SO(3)$ symmetry (angular momentum algebra): it fixes the angular basis, while the radial operator and eigenvalues still depend on the chosen hyperelastic constitutive law and background pre-stress.
\(L^2=-\Delta_{S^2}\), i.e.\ the Laplace--Beltrami operator of the \emph{round} metric on \(S^2\).
This is a symmetry/basis statement about the angular sector and does \emph{not} assume a different
constitutive model for the brane substrate.

Its eigenfunctions are the complex spherical harmonics \(Y_{\ell m}(\theta,\phi)\), whose orthonormality,
completeness on the sphere, \(2\pi\)-periodicity in \(\phi\), and regularity at the poles are automatic
properties of the corresponding self-adjoint eigenproblem.

\subsection{Amplitude scaling and energy normalization}
\label{subsec:amplitude_scaling}

The linear eigenproblem fixes mode \emph{shapes} and frequencies but not overall amplitude:
if \(\eta\) is a mode, so is \(c\,\eta\). Physical amplitudes are selected only when nonlinearities are included
and the total energy is fixed. For an electron-like excitation we impose the rest-energy normalization
\begin{equation}
\overline{E}[\X] \;=\; m_e c^2,
\label{eq:rest_energy_normalization}
\end{equation}
where \(\overline{E}\) denotes the cycle-averaged \emph{excess} energy above the pre-stressed ground state.

In the small-slope regime both kinetic and elastic excess energies scale quadratically with an amplitude
parameter \(A\) (schematically \(\overline{E}\propto A^2\)). The confinement strength depends on how much energy
is effectively coupled into the transverse component, which is controlled by the background pre-stress
(\(\alpha<1\)) and the nonlinear dependence of \(g_{ij}(\X)\) on transverse gradients.


\subsection{Two confinement hypotheses (transverse polarization vs.\ connection-mediated self-trapping)}
\label{subsec:confinement_hypotheses}

The substrate equations \eqref{eq:continuous_eom} permit multiple \emph{interpretations} of what stabilizes an
electron-like localized excitation against dispersion. We state two internally consistent hypotheses that are
distinguishable by numerical experiments. In both cases there is no fundamental point charge; ``charge'' is a
label of the far-field phase/holonomy structure derived from the brane waves.

\paragraph{Hypothesis A (transverse-sector confinement; \(X^4\) as primary binder).}
Localization is dominated by nonlinear geometric coupling into the transverse embedding component \(X^4\)
(controlled by the pre-stress parameter \(\alpha\)). The confined object is a self-consistent nonlinear standing
wave whose core necessarily carries substantial \(X^4\) gradients/polarization, while the emergent U(1) phase and
(optionally) a Wilczek--Zee connection arise as \emph{readouts} of the internal near-degenerate mode subspace
(Sec.~\ref{sec:connections}) and characterize the long-range interaction.

\paragraph{Hypothesis B (connection-mediated self-trapping; gauge structure as the effective binder).}
Localization is dominated by the geometry of the carrier-band eigenbundle: the effective ``gauge field'' is the
Berry/Wilczek--Zee connection induced by adiabatic transport of a near-degenerate polarization subspace
(ideally the two transverse/helicity modes tangent to the physical 3D brane).
The parameter \(\alpha\) controls the degree of band coupling and the splitting of that subspace; when the subspace
is isolated and transport is adiabatic, projection yields a covariant derivative with a rank \(n=1\) U(1) Berry connection
(or \(n>1\) Wilczek--Zee connection). By contrast, strong non-adiabatic mode conversion typically produces beating
and should be treated as explicit inter-band dynamics rather than as curvature of a single Abelian connection.
Microscopically the dynamics is still entirely that of \eqref{eq:continuous_eom}; the ``gauge field'' is not an
independent degree of freedom but a reparametrization of mode-frame transport (Sec.~\ref{sec:connections}).

\paragraph{Operational discriminants (simulation checks).}
\begin{enumerate}
  \item \textbf{Suppress transverse motion.} Constrain \(X^4\equiv 0\) (or heavily penalize \(\nabla X^4\)):
        Hypothesis~A predicts loss of confinement; Hypothesis~B allows partial confinement if the relevant
        internal mixing persists within the remaining components/subspace.
  \item \textbf{Energy partition vs.\ curvature.} Hypothesis~A predicts a persistent core energy fraction in
        \(X^4\)-gradients; Hypothesis~B predicts confinement correlates more strongly with large Berry/Wilczek--Zee
        curvature \((f_{\mu\nu},\bbF_{\mu\nu})\) and with near-degeneracy/branch-mixing measures.
  \item \textbf{\(\alpha\)-dependence.} Vary \(\alpha\) while holding the carrier frequency fixed:
        Hypothesis~A ties confinement primarily to transverse coupling strength; Hypothesis~B ties it primarily to
        the strength of internal branch mixing and the induced connection/holonomy.
\end{enumerate}


\subsection{Discrete degrees of freedom}
We discretize $\Omega$ by a 3D lattice of nodes $p\in\mathcal{V}$.
Each node carries a position $\mathbf{R}_p(t)\in\mathbb{R}^4$ and velocity
$\dot{\mathbf{R}}_p(t)\in\mathbb{R}^4$.

The brane is a graph $G=(\mathcal{V},\mathcal{E})$ with edges $(p,q)\in\mathcal{E}$
representing Hooke springs embedded in $\mathbb{R}^4$.

\subsection{Isometric neighborhoods (rotational symmetry in the continuum limit)}
To approximate isotropic elasticity, the neighborhood stencil must be closed
under the relevant symmetry group.

\paragraph{Preferred option: FCC (high isotropy).}
Use an FCC lattice in material space (nearest-neighbor set with 12 equal-length
directions). This improves isotropy relative to a simple cubic grid.

\paragraph{Fallback option: cubic 26-stencil with distance-based weights.}
On a cubic grid, connect each node to its 26 neighbors with offsets
$\boldsymbol{\delta}\in\{-1,0,1\}^3\setminus\{(0,0,0)\}$, grouped by distance:
\begin{itemize}
\item axial: $|\boldsymbol{\delta}|=1$ (6 neighbors),
\item face-diagonal: $|\boldsymbol{\delta}|=\sqrt{2}$ (12 neighbors),
\item body-diagonal: $|\boldsymbol{\delta}|=\sqrt{3}$ (8 neighbors).
\end{itemize}
To preserve isotropy of the quadratic continuum limit, scale spring stiffness
by the squared distance:
\begin{equation}
  k_{\boldsymbol{\delta}} \;=\; k\,w_{\boldsymbol{\delta}},
  \qquad
  w_{\boldsymbol{\delta}} := \frac{1}{|\boldsymbol{\delta}|^2}
  \;=\;
  \begin{cases}
    1, & |\boldsymbol{\delta}|=1,\\
    1/2, & |\boldsymbol{\delta}|=\sqrt{2},\\
    1/3, & |\boldsymbol{\delta}|=\sqrt{3}.
  \end{cases}
  \label{eq:iso_weights}
\end{equation}
This choice makes the second-moment tensor $\sum_{\delta}w_\delta\,\delta\delta^{\mathsf T}$
proportional to the identity, yielding an isotropic Laplacian in the continuum limit.

\subsection{Discrete stretching energy (rest-length model)}
Let $h_\star$ be the geometric spacing between nearest neighbors (set by the grid
and boundary conditions). For an edge $(p,q)$ associated with offset
$\boldsymbol{\delta}$ we define
\[
  \ell_{pq}(t) := \|\mathbf{R}_q(t)-\mathbf{R}_p(t)\|_{\mathbb{R}^4},
  \qquad
  \ell^0_{pq} := \ell_0\,|\boldsymbol{\delta}|,
\]
where $\ell_0>0$ is the rest length for unit-offset neighbors.

The total potential energy is
\begin{equation}
  U(\{\mathbf{R}_p\})
  = \sum_{(p,q)\in\mathcal{E}} \frac{k_{\boldsymbol{\delta}(p,q)}}{2}\,
    \big(\ell_{pq}-\ell^0_{pq}\big)^2,
  \label{eq:discrete_energy}
\end{equation}
and the kinetic energy is
\begin{equation}
  T(\{\dot{\mathbf{R}}_p\})
  = \sum_{p\in\mathcal{V}} \frac{m}{2}\,|\dot{\mathbf{R}}_p|^2.
  \label{eq:discrete_kinetic}
\end{equation}
The discrete Lagrangian is $L=T-U$.

\subsection{Discrete equations of motion (4D spring forces)}
The force on node $p$ is the sum of spring forces along incident edges:
\begin{equation}
  \mathbf{F}_p
  = -\nabla_{\mathbf{R}_p}U
  = \sum_{q\in\mathcal{N}(p)}
    k_{\boldsymbol{\delta}(p,q)}\,
    \Big(1-\frac{\ell^0_{pq}}{\ell_{pq}}\Big)\,(\mathbf{R}_q-\mathbf{R}_p),
  \label{eq:discrete_force}
\end{equation}
with $\mathcal{N}(p)$ the neighbor set. The dynamics is
\begin{equation}
  m\,\ddot{\mathbf{R}}_p = \mathbf{F}_p + \mathbf{F}^{\mathrm{bc}}_p,
  \label{eq:discrete_eom}
\end{equation}
where $\mathbf{F}^{\mathrm{bc}}_p$ enforces the chosen boundary conditions
(periodic identification or clamping / holding forces).

\subsection{Coupling / pre-stretch parameter (discrete)}
Define the same coupling knob as in the continuum:
\begin{equation}
  \alpha := \frac{\ell_0}{h_\star}\in(0,1],
  \qquad
  \varepsilon := \alpha^{-1}-1.
  \label{eq:discrete_alpha}
\end{equation}
\begin{itemize}
\item $\alpha=1$: rest lengths match geometric spacing (no background tension).
\item $\alpha<1$: uniform pre-tension because the lattice is held at spacing $h_\star$
      while springs prefer $\ell_0$; this amplifies geometric coupling between
      transverse and lateral motion through the nonlinear factor
      $(1-\ell^0_{pq}/\ell_{pq})$ in \eqref{eq:discrete_force}.
\end{itemize}
This is the discrete counterpart of the continuum pre-stretch encoded by the
reference metric scale $\ell_0$ versus the held metric scale $h_\star$.

\subsection{Time integration and numerical stability}
We integrate \eqref{eq:discrete_eom} with a symplectic second-order method (Velocity Verlet).
A conservative stability condition is a CFL-like bound
\begin{equation}
  \Delta t \;\lesssim\; \eta \sqrt{\frac{m}{k}},
  \qquad \eta\in(0,1),
\end{equation}
with the precise constant depending on the stencil and boundary conditions.

\subsection{Discrete-to-continuum parameter correspondence}
With geometric spacing $h_\star$ and node mass $m$, the effective continuum mass
density scales as
\begin{equation}
  \rho_m \sim \frac{m}{h_\star^3}.
\end{equation}
For isotropic stencils, the effective continuum stiffness scales as
\begin{equation}
  K \sim \frac{k}{h_\star},
  \qquad\Rightarrow\qquad
  c^2=\frac{K}{\rho_m}\sim \frac{k\,h_\star^2}{m},
\end{equation}
consistent with the linearized wave speed used for calibration.

% ============================================================
\section{Narrowband Multiscale Ansatz}
\label{sec:multiscale}

\subsection{Carrier scale and slow envelope}
Assume the excitation of interest is narrowband around a fixed angular frequency
\(\omega_0\). In the project, \(\omega_0\) is hardcoded from the Compton wavelength
\(\lambda_C\):
\[
\bar\lambda_C := \frac{\hbar}{m_e c},\qquad \omega_0=\omega_C:=\frac{c}{\bar\lambda_C}=\frac{m_e c^2}{\hbar}.
\]
(If one uses the non-reduced Compton wavelength \(\lambda_C=h/(m_e c)\), then \(\omega_C=2\pi c/\lambda_C\).)
Write the displacement as a slowly modulated carrier:
\begin{equation}
\uvec(\x,t)
\approx
\Re\!\left[
\e^{-\ii \omega_0 t}\sum_{a=1}^{n} \Psi^a(\x,t)\, \mathbf{e}_a(\x,t)
\right],
\label{eq:ansatz}
\end{equation}
where
\begin{itemize}
  \item \(\Psi(\x,t)\in\C^{n}\) is a complex envelope (slow in \(t\) relative to \(\omega_0^{-1}\)),
  \item \(\{\mathbf{e}_a(\x,t)\}_{a=1..r}\subset\R^4\) is a local orthonormal mode frame
        (polarization / internal basis) spanning the retained subspace,
  \item \(n=1\) corresponds to a U(1) (rank-1) bundle; \(n>1\) yields a non-Abelian bundle.
\end{itemize}


\subsection{Ordered narrowband sector as a brane ``order parameter'' (phase rigidity)}
\label{subsec:order_parameter}

Equation~\eqref{eq:ansatz} is more than a convenient signal-processing device: it defines a macroscopic \emph{complex}
envelope field. Write (rank-1 for clarity)
\[
\Psi(\x,t)=A(\x,t)\,\e^{\ii\phi(\x,t)}.
\]
The phase field $\phi$ is meaningful only when the carrier is \emph{coherent}: $A$ must remain non-negligible and the
phase must be correlated over distances long compared to the carrier wavelength. This is the precise sense in which the
brane picture can support an ``order parameter with a phase'' (by analogy to superfluids), without assuming any change
to microscopic statistics: we are describing a \emph{coherent narrowband sector} of a classical substrate, not claiming
that fermionic exclusion is ``turned off''.

\paragraph{Phase rigidity from brane elasticity.}
The origin of phase stiffness is the ordinary gradient energy of the brane.
After averaging over the fast carrier oscillation, the projected envelope energy contains terms of the form
\[
E_{\rm env}\;\sim\;\int \frac{K_{\rm eff}}{2}\,|\nabla\Psi|^2\,\dd^3x \;+\;\cdots
\qquad\Rightarrow\qquad
|\nabla\Psi|^2 = |\nabla A|^2 + A^2|\nabla\phi|^2.
\]
Thus spatial phase twists cost energy whenever $A$ is nonzero. This ``rigidity'' is what distinguishes an ordered
narrowband sector from a broadband vibration bath.

\paragraph{Why a ``vibration soup'' does not define a gauge field.}
A broadband superposition with essentially random phases has no stable envelope $\Psi$ and no long-lived $\phi(\x,t)$.
In that regime, mode conversion and dephasing make the measured frame transport noisy; the Berry connection becomes a
short-correlation diagnostic rather than a clean emergent gauge potential.
Practically, the U(1) reduction requires (i) a narrowband carrier, (ii) a large phase-coherence length, and
(iii) an adiabatically varying mode frame (Sec.~\ref{sec:connections}).

\paragraph{Role of $\alpha$ and branch structure.}
When $\alpha$ produces strong longitudinal--transverse separation and little mixing, a Berry description based on
\emph{longitudinal/transverse hybridization} is suppressed because the eigenmodes are already close to pure branches.
Conversely, at intermediate $\alpha$ (or near avoided crossings) the instantaneous eigenmodes are hybrid and the
associated polarization frame can vary with propagation direction and background deformation, enabling nontrivial
Berry holonomy. Even in the decoupled limit, a U(1) (or non-Abelian) connection may still arise from transport within
a near-degenerate \emph{transverse polarization} subspace; the present framework treats both cases uniformly by
defining $\mathbf{E}$ as the appropriate eigenframe of the linearized operator.

\subsection{Complex quadrature from real simulation variables}
In a real-valued simulator, construct the complex amplitude from displacement and velocity
(\texttt{branesim.berry.analysis}):
\begin{equation}
\mathbf{a}(\x,t) = \sqrt{\omega_0}\,\uvec(\x,t) + \ii\,\frac{\vvec(\x,t)}{\sqrt{\omega_0}}
\in \C^{E},
\label{eq:complex_amp}
\end{equation}
where \(E\in\{2,3,4\}\) is the number of embedding components used for the analysis.
Normalize to a ray:
\[
\ket{\psi(\x,t)} = \frac{\mathbf{a}(\x,t)}{\|\mathbf{a}(\x,t)\|}.
\]
Undefined regions are masked by amplitude thresholds.

\subsection{What the Berry phase is ``of'': eigenmode bundles, not arbitrary decompositions}
The mode-frame $\mathbf{E}(x)=[\mathbf{e}_1,\dots,\mathbf{e}_n]$ must be tied to the \emph{instantaneous eigenmodes}
of the \emph{linearized} brane operator around a chosen background $\X^\star(x)$ and in the chosen carrier band near $\omega_0$.
Concretely, let $\mathcal{L}_{\X^\star}$ denote the quadratic (small-amplitude) operator governing carrier-scale dynamics.
For each wavevector (or local ray direction) and band index $b$, the physical polarization vectors are eigenvectors
$\mathbf{e}_b$ of $\mathcal{L}_{\X^\star}$; the Berry/Wilczek--Zee connection is the connection on the resulting eigenbundle.
This matters whenever ``longitudinal'' and ``transverse'' components \emph{mix}: in that case $L/T$ are \emph{not} eigenmodes,
and the correct $\mathbf{E}$ is a \emph{hybrid} eigenframe (bands $b=\pm$), not a kinematic $L/T$ projection.

\subsection{Adiabatic single-band regime vs.\ beating (mode conversion)}
A \emph{clean} rank-1 (U(1)) Berry description requires that the dynamics remains in a \emph{single isolated band}:
inter-band transitions must be negligible on the envelope scale. If the initial condition excites a superposition of two nearby
bands, the resulting amplitude modulation (``beating'') is a \emph{dynamical interference} effect, not geometric phase.
In practice, the adiabatic/single-band condition is enforced by a spectral separation
$|\omega_b-\omega_{b'}|\gg \|\partial_t\|$ together with slow spatial variation of the eigenframe.

\paragraph{Dephasing from broadband backgrounds.}
If the brane hosts a broadband background of random vibrations, the envelope-phase $\phi$ and the eigenframe $\mathbf{E}$
typically decorrelate on short scales, effectively reintroducing inter-band transitions and making $a_\mu$ noisy.
In this regime a measured ``Berry signal'' should be interpreted as a diagnostic of mode conversion rather than as an
effective gauge potential. This motivates the explicit coherence assumption in Sec.~\ref{sec:maxwell_term}.


\subsection{Propagation speed and (possible) dispersion of the emergent gauge sector}
The causal propagation of features in the measured connection $a_\mu$ is inherited from the physical channel that changes the
eigenframe (transport of a narrowband packet and/or propagation of background deformations). In the single-band regime this
is governed by the \emph{group velocity} $v_g(\omega)=\dd\omega/\dd k$ of the chosen band in the operating band around $\omega_0$.
Thus, frequency dependence enters only through \emph{dispersion}: on the discrete lattice and near avoided crossings induced by
band coupling, $v_g$ can vary with $\omega$; in the long-wavelength continuum regime $v_g$ is approximately constant.


% ============================================================
\section{Connections from Mode-Frame Transport (Berry / Wilczek--Zee)}
\label{sec:connections}

\subsection{Rank-1 (U(1)) Berry connection}
For \(n=1\), let \(\ket{u(x)}\) be the normalized local mode, with \(x^\mu=(c_{\rm EM}t,\x)\).
Define
\[
a_\mu(x) = \ii \braket{u(x)}{\partial_\mu u(x)}.
\]
The Berry phase along a path \(\Gamma\) is \(\gamma_B=\int_\Gamma a_\mu\,\dd x^\mu\),
and the curvature is
\[
f_{\mu\nu} = \partial_\mu a_\nu - \partial_\nu a_\mu.
\]


\paragraph{Geometric vs.\ dynamical phase (important distinction).}
The connection above is geometric: it is defined by the transport of an \emph{eigenmode} (or eigen-subspace).
A superposition of nearby bands can produce beating patterns through $\Delta\omega=\omega_+-\omega_-$, but such beating is not a
Berry curvature signal and generally indicates non-adiabatic mode conversion. The Berry/U(1) reduction is therefore tied to
(i) an isolated band (or a symmetry-protected helicity sector) and (ii) slow variation of the associated eigenframe.


\subsection{Non-Abelian (Wilczek--Zee) connection for internal electron subspace}
For \(n>1\), choose an orthonormal basis \(\ket{u_a(x)}\), \(a=1..r\), and define
\[
\bbA_\mu^{ab}(x) = \ii \braket{u_a(x)}{\partial_\mu u_b(x)}\in\mathfrak{u}(n).
\]
Curvature:
\[
\bbF_{\mu\nu} = \partial_\mu \bbA_\nu - \partial_\nu \bbA_\mu - \ii[\bbA_\mu,\bbA_\nu].
\]

\subsection{Envelope covariant derivative}
Inserting the ansatz \eqref{eq:ansatz} and projecting onto the retained subspace produces,
generically, a covariant derivative acting on \(\Psi\):
\begin{equation}
\covD_\mu \Psi
=
\left(\partial_\mu + \ii a_\mu \mathbf{1}_n + \ii\bbA_\mu\right)\Psi.
\label{eq:covD}
\end{equation}

\paragraph{Derivation sketch.}
\begin{enumerate}
  \item Write \(\uvec \approx \Re(\e^{-\ii\omega_0 t} \mathbf{E}\Psi)\) where
        \(\mathbf{E}(x)=[\mathbf{e}_1,\dots,\mathbf{e}_r]\in\R^{4\times r}\).
  \item Differentiate and project with \(\mathbf{E}^\top\) (or \(\mathbf{E}^\dagger\) in a complexified
        inner product) to obtain terms of the form \((\mathbf{E}^\top \partial_\mu \mathbf{E})\Psi\).
  \item Identify the antisymmetric part as the gauge connection; show basis changes
        \(\mathbf{E}\to \mathbf{E}U(x)\) induce the standard transformation law for \(\bbA_\mu\).
\end{enumerate}

% ============================================================
\section{QED-like Effective Lagrangian (Kinematics)}
\label{sec:qed_like}

\subsection{Local phase symmetry}
Let \(\psi(x)\) be a Dirac spinor envelope field. Under a local U(1) phase change
\(\chi(x)\):
\[
\psi \to \e^{\ii \frac{q}{\hbar}\chi(x)}\psi,
\qquad
A_\mu \to A_\mu - \partial_\mu \chi.
\]

\subsection{Meaning of the U(1) phase in this framework}
In standard QED the function $\chi(x)$ expresses a redundancy of description (a choice of local phase convention).
In the present brane framework the same redundancy is realized as a choice of local mode-frame
$\mathbf{E}(x)\to \mathbf{E}(x)U(x)$, and the emergent potential is the corresponding Berry connection.
Observable quantities depend only on the gauge-invariant curvature $F_{\mu\nu}$, not on $\chi(x)$ itself.




\paragraph{Envelope phase vs.\ gauge choice.}
The overall envelope phase $\phi$ in $\Psi=A\e^{\ii\phi}$ is a physical descriptor of the ordered narrowband sector
(Sec.~\ref{subsec:order_parameter}) and inherits an elastic ``phase stiffness'' from gradient energy.
The gauge redundancy $\chi(x)$, by contrast, expresses the freedom to rotate the local mode-frame
$\mathbf{E}(x)\to \mathbf{E}(x)U(x)$ used to represent that sector. The Berry connection $a_\mu$ records how this
frame twists in spacetime; only its curvature is frame independent.

\subsection{Minimal coupling with Berry-induced potential}
Identify the physical EM potential with the Berry connection up to units:
\[
A_\mu := \frac{\hbar}{q}\,a_\mu
\quad\Rightarrow\quad
\partial_\mu + \ii \frac{q}{\hbar}A_\mu = \partial_\mu + \ii a_\mu.
\]

\subsection{Compact (2-term) and expanded (3-term) forms}
\paragraph{2-term (compact).}
\begin{equation}
\Lag
=
\psibar\left(\ii\hbar c\,\gamm^\mu \covD_\mu - mc^2\right)\psi
-
\frac{1}{4\mu_0}F_{\mu\nu}F^{\mu\nu},
\label{eq:Lag_compact}
\end{equation}
with \(\covD_\mu=\partial_\mu + \ii\frac{q}{\hbar}A_\mu\).

\paragraph{3-term (expanded).}
\begin{equation}
\Lag
=
\underbrace{\psibar(\ii\hbar c\,\gamm^\mu \partial_\mu - mc^2)\psi}_{\text{Dirac}}
+
\underbrace{(-q\,\psibar\gamm^\mu\psi)A_\mu}_{\text{coupling}}
-
\underbrace{\frac{1}{4\mu_0}F_{\mu\nu}F^{\mu\nu}}_{\text{EM field}},
\label{eq:Lag_expanded}
\end{equation}

\subsection{Including the Wilczek--Zee sector}
If the electron envelope carries an internal index (dimension \(r\)), replace \(\psi\) with
\(\Psi\in\C^{4}\otimes\C^{n}\) (Dirac spinor valued in the internal subspace) and use
\[
\covD_\mu \Psi = \left(\partial_\mu + \ii a_\mu \mathbf{1}_n + \ii\bbA_\mu\right)\Psi.
\]
This is the place to state explicitly what \(r\) represents in the brane model
(e.g., internal polarization/standing-wave basis supporting spin-like behavior).

% ============================================================

\section{Dynamics: Emergent Maxwell Term from Coarse-Graining}
\label{sec:maxwell_term}

Sections~\ref{sec:connections} and~\eqref{eq:covD} provide \emph{gauge-covariant kinematics}:
a slowly varying mode-frame produces a connection \(a_\mu\) (or \(\bbA_\mu\)) and
the envelope dynamics depends on it through the covariant derivative.
Electromagnetism, however, is not only kinematics. It requires an autonomous field equation
for the connection---equivalently, an effective action for \(a_\mu\) whose Euler--Lagrange
equations are Maxwell-like.

This section provides a ``paper-grade'' derivation at the level of an \emph{effective field theory}
(EFT): we state the approximation regime precisely, perform the standard multi-scale reduction,
and show that the \emph{leading} local, gauge-invariant, quadratic term in the resulting effective
action is necessarily of Maxwell/Yang--Mills form. The only part not closed analytically is the
numerical value of the prefactor, which is a material/geometry constant of the brane and can be
measured in simulation.

\subsection{Assumptions (explicit approximation regime)}
We work in a regime where the following hold.
\begin{enumerate}
  \item \textbf{Separation of scales.} There exists a carrier frequency \(\omega_0\) (Compton-scale in our calibration)
        and slow envelope variations such that \(\|\partial_\mu \Psi\|\ll \omega_0\|\Psi\|\) and
        \(\|\partial_\mu \mathbf{E}\|\ll \omega_0\|\mathbf{E}\|\), where \(\mathbf{E}(x)=[\mathbf{e}_1,\dots,\mathbf{e}_n]\)
        is an orthonormal local mode-frame spanning the relevant \(n\)-dimensional internal subspace.

  \item \textbf{Coherence (ordered narrowband sector).} Random broadband vibrations are subdominant on the carrier scale:
        the complex envelope has a long phase-coherence length and time (relative to the carrier wavelength/period), and the
        mode-frame transport extracted from $\mathbf{E}(x)$ is smooth after coarse-graining. In practice this corresponds to
        maintaining a narrow spectral bandwidth around $\omega_0$ and suppressing non-adiabatic band hopping/dephasing.
  \item \textbf{Single-band adiabaticity.} The retained mode-frame $\mathbf{E}(x)$ spans an \emph{isolated} carrier band (rank 1) or a \emph{well-separated} rank-$n$ subspace.
        Inter-band transitions (mode conversion) are negligible on the envelope scale; otherwise beating and non-adiabatic mixing invalidate an Abelian U(1) reduction and require an explicit multi-band (Wilczek--Zee) treatment.
  \item \textbf{Near-linearity at the carrier scale.} Around the chosen background configuration \(\X^\star\),
        the carrier-scale dynamics is well-approximated by the quadratic expansion of the brane action.
        Nonlinearities remain essential for confinement (electron sector) but appear as slow, envelope-scale effects
        in the present reduction.
  \item \textbf{Locality, isotropy, and the meaning of \(c\) in the EFT.}
        After averaging over a few carrier periods and a few lattice spacings, the resulting EFT is local and
        (approximately) Lorentz-invariant in the \emph{long-wavelength} regime of the selected EM subbundle.
        The characteristic causal speed entering \(x^\mu=(c_{\rm EM}t,\x)\) is the group velocity
        \(c_{\rm EM}\approx v_g(\omega_0)\) of that band in the operating narrowband.
        In the reference implementation we calibrate parameters such that \(c_{\rm EM}\) matches the physical \(c\)
        for low \(k\). Outside this regime (near lattice scales or near strong band coupling) dispersion produces
        higher-derivative corrections and small Lorentz-violating effects.

  \item \textbf{Gauge redundancy = frame choice.} The frame \(\mathbf{E}\) is only defined up to a local rotation
        \(\mathbf{E}(x)\to \mathbf{E}(x)\,U(x)\) with \(U(x)\in U(1)\) (or \(U(n)\) in the Wilczek--Zee case),
        and physical predictions may not depend on this choice.
  \item \textbf{No explicit symmetry breaking in the EM sector.} There is no term in the coarse-grained energy that fixes a preferred
        absolute phase (rank-1) or absolute internal orientation (rank-\(n\)) of the EM mode-frame. (If such a term existed, it would generate
        a Proca-like mass; see below.)
\end{enumerate}

\subsection{Multi-scale reduction and covariant envelope terms (recap)}
Start from the brane action
\[
S[\X]=\int \dd t \int_{\Omega}\!\left[\frac{\rho}{2}\,|\partial_t\X|^2 - W\!\big(g(\X);g^0\big)\right]\dd^3x,
\qquad g_{ij}=\partial_i\X\cdot\partial_j\X.
\]
Linearizing around a background \(\X^\star\) yields a second-order operator \(\mathcal{L}\) such that (schematically)
\(\rho\,\partial_{tt}\uvec + \mathcal{L}\uvec = 0\), where \(\uvec\) is the small displacement.
In the narrowband regime we use the ansatz (cf.~Sec.~\ref{sec:connections})
\begin{equation}
\uvec(x,t)\;\approx\;\Re\!\left[\e^{-\ii\omega_0 t}\,\mathbf{E}(x)\,\Psi(x)\right],
\qquad
\mathbf{E}^\top\mathbf{E}=\mathbb{I}_n,
\label{eq:narrowband_ansatz}
\end{equation}
where \(\Psi(x)\in\C^n\) is the complex envelope in the chosen mode-frame.
Projecting the averaged Lagrangian onto the span of \(\mathbf{E}\) and keeping only leading terms in the slow-variation expansion gives
\begin{equation}
S_{\rm env}[\Psi;\bbA]\;=\;\int \dd^4x\;\Big(
\frac{\kappa}{2}\,(D_\mu\Psi)^\dagger (D^\mu\Psi) - V_{\rm eff}(\Psi)
+\cdots\Big),
\label{eq:Seff_env}
\end{equation}
with covariant derivative
\(
D_\mu\Psi = (\partial_\mu + \ii\,\bbA_\mu)\Psi
\)
and connection
\(
\bbA_\mu = \ii\,\mathbf{E}^\dagger\partial_\mu \mathbf{E}
\)
(up to the inner-product convention used for the projected complexified subspace).
The dots denote higher-derivative corrections suppressed by \(\omega_0^{-1}\).
Equation~\eqref{eq:Seff_env} is the \emph{matter} part of the effective theory. We now derive the \emph{field} part.

\subsection{From brane energy to a curvature penalty}
The key point is that \(\bbA_\mu\) is not an independent fundamental field: it is a collective coordinate describing
how the local internal mode-frame twists in spacetime. Any coarse-grained energy associated with such twisting must respect
the gauge redundancy \(\mathbf{E}\to \mathbf{E}U(x)\).
Consequently, the effective action for \(\bbA_\mu\) can only be built from gauge-invariant tensors, i.e.
the curvature
\begin{equation}
\bbF_{\mu\nu} = \partial_\mu \bbA_\nu - \partial_\nu \bbA_\mu - \ii[\bbA_\mu,\bbA_\nu]
\qquad (\,\text{rank }n\,),
\label{eq:curvature_def}
\end{equation}
or \(f_{\mu\nu}=\partial_\mu a_\nu-\partial_\nu a_\mu\) in the rank-1 case.

Under the locality + derivative-expansion assumptions, the leading term in the pure-connection EFT is the lowest-order
scalar that is (i) gauge invariant, (ii) parity even, and (iii) quadratic in fields and derivatives. Up to total derivatives,
this is uniquely
\begin{equation}
S_{\rm gauge}[\bbA]\;=\;-\frac{1}{2g_{\rm eff}^2}\int \dd^4x\;\tr\!\left(\bbF_{\mu\nu}\bbF^{\mu\nu}\right)
\quad\text{(rank \(n\))},
\label{eq:Seff_YM}
\end{equation}
reducing to the Maxwell form for \(n=1\),
\begin{equation}
S_{\rm gauge}[a]\;=\;-\frac{1}{4\mu_0^{\rm eff}}\int \dd^4x\; f_{\mu\nu}f^{\mu\nu}.
\label{eq:Seff_Maxwell}
\end{equation}
The prefactors \(g_{\rm eff}\) or \(\mu_0^{\rm eff}\) are not free constants in the full model: they are determined by the
elastic response of the brane and the details of the coarse-graining. In practice, they can be obtained by a numerical ``calibration experiment''
(see Sec.~\ref{sec:maxwell_calibration}).

\paragraph{Why the Maxwell/Yang--Mills form is inevitable in this regime.}
Any local quadratic action for a connection has the schematic structure
\(\int \bbA_\mu\,\Pi^{\mu\nu}(\partial)\,\bbA_\nu\).
Gauge invariance under \(\bbA_\mu\to U^{-1}\bbA_\mu U + \ii U^{-1}\partial_\mu U\) forces \(\Pi^{\mu\nu}\) to be transverse,
i.e. proportional to \(\eta^{\mu\nu}\partial^2-\partial^\mu\partial^\nu\) at leading derivative order.
This transversality is exactly what is encoded by \(\tr(\bbF_{\mu\nu}\bbF^{\mu\nu})\).
Thus, once (i) gauge redundancy and (ii) locality are accepted, the Maxwell/Yang--Mills term is the universal leading operator.

\subsection{Field equations and sources}
Varying \(S_{\rm gauge}+S_{\rm env}\) with respect to the connection yields Maxwell/Yang--Mills equations with a source term
coming from the envelope sector:
\begin{align}
\partial_\mu f^{\mu\nu} &= \mu_0^{\rm eff}\, j^\nu
\qquad\text{(rank \(1\))},\\
D_\mu \bbF^{\mu\nu} &= g_{\rm eff}^2\, \mathbf{J}^\nu
\qquad\text{(rank \(n\))},
\end{align}
where \(j^\nu\) (or \(\mathbf{J}^\nu\)) is the Noether current associated with the gauge redundancy of the projected envelope action.
In the rank-1 case, to compare with the standard QED normalization one may identify the physical potential by a scale factor
\(A_\mu = (\hbar/q)\,a_\mu\), which turns the Berry curvature into the usual field tensor \(F_{\mu\nu}=\partial_\mu A_\nu-\partial_\nu A_\mu\)
and makes the coupling constant appear in the matter term rather than in the definition of the connection (cf.~Sec.~\ref{sec:qed_like}).

\subsection{Why the emergent gauge field is massless here}
A mass term would have the form \(\int m_A^2\,\tr(\bbA_\mu\bbA^\mu)\) (Proca).
Such a term is \emph{not} gauge invariant and therefore cannot appear in the EFT unless the gauge redundancy is explicitly or spontaneously broken.
Under Assumption~(5), the first allowed pure-connection operator is \(\tr(\bbF_{\mu\nu}\bbF^{\mu\nu})\), which describes a massless gauge field
with two propagating transverse polarizations (after gauge fixing and constraints).
This provides the precise sense in which ``the Berry gauge field is automatically massless'' in the present framework:
masslessness is protected by the frame-gauge redundancy rather than by an additional postulate.

\subsection{Polarization count and suppression of extra embedding modes}
In 3+1 dimensions, a massless \(U(1)\) gauge field has two physical polarizations.
The remaining components are removed by gauge choice and the Gauss constraint.
In our brane picture this corresponds to selecting the \emph{EM subbundle} as a two-dimensional transverse polarization plane
tangent to the physical 3D brane, and treating other embedding-direction excitations as belonging to different sectors:
\begin{itemize}
  \item \textbf{Longitudinal/timelike components} of \(a_\mu\) do not represent additional propagating modes; they enforce constraints
        (e.g. Gauss law) once a gauge is fixed.
  \item \textbf{Out-of-plane polarization (fourth embedding component).} Excitations predominantly in the fourth embedding direction
        are not assigned to the EM subbundle. They can be (i) energetically gapped by the background configuration and boundary conditions,
        (ii) localized as part of confined particle modes, or (iii) suppressed by the choice of projection defining \(\mathbf{E}\).
        Which mechanism dominates is a quantitative question addressed by simulation diagnostics.
\end{itemize}

\subsection{Calibration in simulation: extracting the effective permeability from energy vs curvature}
\label{sec:maxwell_calibration}
Equation~\eqref{eq:Seff_Maxwell} is predictive once its prefactor is fixed.
A direct numerical way to obtain \(\mu_0^{\rm eff}\) is:
\begin{enumerate}
  \item Prepare a configuration with a controlled, slowly varying twist of the local EM mode-frame (e.g. by boundary driving or by constructing
        a long-wavelength packet in the EM subbundle).
  \item Compute the coarse-grained curvature \(f_{\mu\nu}\) (or \(\bbF_{\mu\nu}\)) from the measured mode-frame transport
        using the discrete analog of~\eqref{eq:curvature_def}.
  \item Measure the corresponding coarse-grained brane energy increment \(\Delta E\) relative to a reference configuration with \(f_{\mu\nu}\approx 0\).
  \item Fit \(\Delta E\) to \(\frac{1}{4\mu_0^{\rm eff}}\int f_{\mu\nu}f^{\mu\nu}\,\dd^3x\) (for static configurations) or the full spacetime action
        (for propagating pulses). The same experiment also checks that the propagation speed of small EM packets matches \(c\) within numerical error.
\end{enumerate}

% ============================================================
\section{Simulation Protocol (Reference Implementation)}
\label{sec:simulation}

\subsection{Electron initialization (rotating spherical-harmonic cavity seed)}
\label{subsec:electron_init}

The reference implementation initializes an ``electron'' as a localized, rotating standing-wave pattern on the 3D brane,
embedded in \(\R^4\). The initializer is designed to (i) create a controlled angular phase winding
(\(m\neq 0\)) in the intrinsic \((x,y,z)\) coordinates, (ii) realize that complex phase rotation in a
\emph{real-valued} simulator by using two quadratures, and (iii) ensure immediate coupling between the
spatial embedding components \((X^1,X^2,X^3)\) and the transverse embedding component \(X^4\) through
the 4D rest-length spring geometry.

\paragraph{Reference-space spherical coordinates.}
All mode functions are evaluated in the intrinsic/reference coordinates \(\x\in\Omega\subset\R^3\).
Given a chosen center \(\x_c\in\R^3\), define \(\mathbf{r}=\x-\x_c\), \(r=\|\mathbf{r}\|\), and
\[
\theta=\arccos\!\left(\frac{z}{r}\right),\qquad
\phi=\operatorname{atan2}(y,x),
\]
with the usual inverse mapping
\(
x=r\sin\theta\cos\phi,\;
y=r\sin\theta\sin\phi,\;
z=r\cos\theta.
\)

\paragraph{Angular factor \(Y_{\ell m}\).}
The angular structure is given by the spherical harmonic \(Y_{\ell m}(\theta,\phi)\), with standard
properties:
\begin{align}
\int_{S^2} Y_{\ell m}(\Omega)\,Y_{\ell' m'}^*(\Omega)\,d\Omega &= \delta_{\ell\ell'}\delta_{mm'},\\
\sum_{\ell=0}^{\infty}\sum_{m=-\ell}^{\ell} Y_{\ell m}(\Omega)\,Y_{\ell m}^*(\Omega')
&= \delta(\Omega-\Omega'),\\
Y_{\ell m}(\theta,\phi+2\pi) &= Y_{\ell m}(\theta,\phi),
\end{align}
and \(Y_{\ell m}\) is bounded and single-valued on \(S^2\) for integer \(\ell\ge 0\), \(|m|\le \ell\).
Choosing \(m\neq 0\) creates a nontrivial azimuthal phase winding \(\sim e^{im\phi}\), which is the
seed for a circulating phase gradient around the symmetry axis.
In the continuum theory, these properties are not imposed by modifying the equations; they arise automatically
because the linearized operator about a spherically symmetric background is self-adjoint and rotation-invariant.

\paragraph{Radial factor \(R_{\ell n}\).}
We use a spherical-Bessel standing-wave profile as a \emph{convenient seed basis} (exact for a scalar Laplacian with Dirichlet boundary conditions); the true linear eigenfunctions in the hyperelastic model are obtained by solving the corresponding self-adjoint radial eigenproblem numerically.
\[
R_{\ell n}(r)= j_\ell(k_{\ell n} r)\,w(r),
\qquad
k_{\ell n}=\frac{\alpha_{\ell n}}{a},
\]
where \(j_\ell\) is the spherical Bessel function, \(\alpha_{\ell n}\) is the \(n\)-th positive root of
\(j_\ell(\alpha)=0\), and \(a\) is the chosen electron radius (cavity radius) in intrinsic units.
The window \(w(r)\) is a smooth taper to avoid a hard discontinuity at \(r=a\); in code we use a
\(\tanh\)-ramp,
\[
w(r)=\frac12\left[1-\tanh\!\left(\frac{r-a}{a/s}\right)\right],
\]
with softness parameter \(s>0\) (larger \(s\) yields a softer edge).

\paragraph{Complex mode and real quadratures.}
Define the complex scalar mode
\[
\psi(\x)=R_{\ell n}(r)\,Y_{\ell m}(\theta,\phi),
\qquad
\psi(\x,t)=\psi(\x)\,e^{-i\omega t}.
\]
A real-valued integrator represents the rotating phase by embedding \(\Re\psi\) and \(\Im\psi\) into a
2D polarization plane in \(\R^4\). Let \(\mathbf{p}_1,\mathbf{p}_2\in\R^4\) be orthonormal vectors spanning this
plane (\(\mathbf{p}_i\cdot\mathbf{p}_j=\delta_{ij}\)). The initializer sets
\begin{align}
\uvec(\x,0) &= A\big(\Re\psi(\x)\,\mathbf{p}_1+\Im\psi(\x)\,\mathbf{p}_2\big),\\
\vvec(\x,0) &= A\,\omega\big(\Im\psi(\x)\,\mathbf{p}_1-\Re\psi(\x)\,\mathbf{p}_2\big),
\end{align}
which corresponds to a local circular rotation in the \((\mathbf{p}_1,\mathbf{p}_2)\) plane at angular frequency
\(\omega\). We choose
\[
\omega = c_{\mathrm{wave}}\,k_{\ell n},
\]
with \(c_{\mathrm{wave}}\) the brane wave speed in the (dimensionless) simulation calibration
(\(c_{\mathrm{wave}}=1\) in the current unit convention).

\paragraph{Amplitude normalization.}
The raw displacement field is scaled such that the maximum Euclidean norm of \(\uvec(\x,0)\) equals the
requested amplitude \(A\), i.e. \(A=\max_{\x}\|\uvec(\x,0)\|\) after normalization.

\paragraph{Polarization choices and \(X^4\) coupling.}
The code supports several fixed polarization presets (e.g. \texttt{xy}, \texttt{xz}, \texttt{yz},
\texttt{spatial}, \texttt{all}) and also accepts user-specified \(\mathbf{p}_1,\mathbf{p}_2\).
For electron seeding we typically use \texttt{spatial\_x4}, i.e. a plane that has support in the spatial
coordinates \emph{and} in \(X^4\). This ensures that \(X^4\) carries oscillatory content and momentum at
\(t=0\) and couples immediately to \((X^1,X^2,X^3)\) through the 4D rest-length spring geometry.

\paragraph{Initializer parameter set (implementation-facing).}
The electron seed is fully specified by the following parameters:
\begin{itemize}
  \item \(\ell,m,n\): spherical-harmonic degree/order and radial root index.
  \item \(a\): electron radius (intrinsic cavity radius) used in \(k_{\ell n}=\alpha_{\ell n}/a\).
  \item \(\x_c\): electron center in intrinsic coordinates.
  \item \(A\): target displacement amplitude after normalization.
  \item \(c_{\mathrm{wave}}\): wave speed used to set \(\omega=c_{\mathrm{wave}}k_{\ell n}\) (unity in current calibration).
  \item Polarization choice: either a preset plane (e.g. \texttt{spatial\_x4}) or explicit orthonormal vectors
        \(\mathbf{p}_1,\mathbf{p}_2\in\R^4\).
  \item \(s\): taper softness used in the window \(w(r)\) at the edge \(r\approx a\).
  \item Flag \texttt{set\_velocities}: whether to set initial velocities from the rotating quadrature formula.
  \item Flag \texttt{normalize}: whether to rescale the seed to match the requested amplitude.
\end{itemize}

\paragraph{Design intent and current limitations.}
This initialization provides a reproducible, parameterized starting configuration with a controllable
azimuthal winding number \(m\) and a well-defined rotating phase in a real-valued simulation.
A long-lived confined ``electron'' would correspond to a self-consistent nonlinear eigenmode of the
coupled 4D rest-length elasticity system. The initializer alone does not enforce that equilibrium; it
provides the candidate initial condition from which such equilibria can be sought and tested.



\subsection{Dispersion, band identification, and adiabaticity checks}
A Berry/U(1) interpretation is only meaningful in an operating regime where a single carrier band can be followed adiabatically.
We therefore recommend the following diagnostics in the reference implementation:
\begin{enumerate}
  \item \textbf{Dispersion scan.} Excite near-plane waves at several wavelengths \(\lambda\) (or wave numbers \(k\)) in the
        candidate EM subbundle and measure \(\omega(k)\) by temporal FFT at probes. Compute the group velocity
        \(v_g(k)=\dd\omega/\dd k\). A nearly constant \(v_g\) indicates an effectively non-dispersive operating band.
  \item \textbf{Mode purity / branch leakage.} For each excited \(k\), project the displacement (or complex amplitude) onto the
        instantaneous eigenframe \(\mathbf{E}\) of the linearized operator. Track power leakage into orthogonal bands.
        Pronounced periodic exchange (beating) indicates a multi-band superposition and invalidates a rank-1 reduction.
  \item \textbf{High-frequency suppression mechanism.} If leakage decreases at higher \(\omega\), distinguish the causes:
        (i) phase-mismatch \(\Delta k(\omega)=|k_+-k_-|\) averaging out conversion, (ii) transverse band edges / cutoff on the lattice,
        or (iii) a genuinely small coupling \(g(k)\) in that band.
\end{enumerate}

\paragraph{Ordered narrowband coherence (``condensate-like'' sector).}
The U(1) reduction presumes that the simulated brane contains a \emph{coherent narrowband carrier}
whose complex envelope \(\Psi=Ae^{i\phi}\) is smooth on scales large compared to the carrier wavelength.
In practice, verify this directly from probe time-series:
\begin{enumerate}
  \item \textbf{Carrier frequency and bandwidth.} At a set of probes \(\{\x_p\}\), form a complex signal \(z_p(t)\).
        Use either (i) the two-quadrature representation already present in the initializer, or (ii) an analytic signal
        \(z=u+i\,\mathcal{H}[u]\) (Hilbert transform). Let \(\omega_0\) be the dominant spectral peak of \(z\) and
        \(\Delta\omega\) its half-power width. Require a narrowband regime, e.g.\ \(\Delta\omega/\omega_0\ll 1\),
        and track \(\Delta\omega(t)\) to detect dephasing.
  \item \textbf{Instantaneous phase diffusion.} Define \(\phi_p(t)=\arg z_p(t)\) and \(\omega_{\rm inst}(t)=\dot{\phi}_p(t)\).
        Remove the mean drift \(\omega_0 t\) and estimate phase diffusion via
        \(D_\phi(\tau)=\langle(\phi_p(t+\tau)-\phi_p(t)-\omega_0\tau)^2\rangle_t/(2\tau)\).
        A stable ordered sector has small \(D_\phi\) over many carrier periods.
  \item \textbf{Spatial phase correlation length.} Define a phase-correlation function
        \(C_\phi(\mathbf{r})=\langle e^{i(\phi(\x+\mathbf{r})-\phi(\x))}\rangle_{\x}\) (and time-average over a steady window).
        Extract a coherence length \(L_\phi\) from \(|C_\phi(r)|\sim e^{-r/L_\phi}\) or the first \(e^{-1}\) crossing.
        A ``vibration soup'' corresponds to short \(L_\phi\) (few wavelengths) and rapidly varying \(\phi\).
  \item \textbf{Adiabaticity against band crossings.} Let \(\Delta\omega_{\rm gap}\) be the frequency gap between the tracked band
        and the nearest orthogonal band at the same \(k\). Estimate the mode-frame rotation rate (from projections or local eigenframes)
        \(\Omega_{\rm mix}\sim \|\chi^\dagger \dot{\chi}\|\) (or its spatial analogue \(\|\chi^\dagger \nabla\chi\|\,v_g\)).
        The rank-1 Berry reduction requires \(\Omega_{\rm mix}\ll \Delta\omega_{\rm gap}\).
\end{enumerate}

\paragraph{Where to compute (reference code path).}
Instrument the solver to record probe time-series \(u(\x_p,t)\) (and, if available, the quadrature pair) at a fixed cadence.
In the Python test harness, store these arrays (e.g.\ \texttt{.npz}) and a summary table (CSV) in the run folder using the
\texttt{TestRunManager.get\_data\_path(...)} helpers (cf.\ \texttt{tests/test\_run\_manager.py}).
A post-analysis script (recommended: \texttt{tests/diagnostics\_coherence.py}) should then:
(i) FFT each probe to extract \(\omega_0,\Delta\omega\),
(ii) compute \(D_\phi(\tau)\) from unwrapped phase,
(iii) compute \(C_\phi(r)\) on a probe lattice (or radial lines),
and (iv) combine these with the branch-leakage projections above into a single
\texttt{coherence\_metrics.csv} plus plots saved via \texttt{TestRunManager.get\_plot\_path(...)}.

\subsection{Berry/Wilczek--Zee extraction and ``U(1)-ness''}
Given a measured mode-frame \(\mathbf{E}(x)\), compute the discrete connection and curvature (as already implemented in
\texttt{branesim.berry.analysis}) and report gauge-invariant diagnostics:
\begin{enumerate}
  \item \textbf{Curvature energy proxy.} Compare \(\int f_{\mu\nu}f^{\mu\nu}\) (rank 1) or \(\int \mathrm{tr}(\bbF_{\mu\nu}\bbF^{\mu\nu})\)
        (rank \(n\)) against the measured coarse-grained brane energy excess to calibrate \(\mu_0^{\rm eff}\)
        (Sec.~\ref{sec:maxwell_calibration}).
  \item \textbf{Abelian reduction check.} In a purported U(1) regime, verify that off-diagonal components are negligible
        (for \(n>1\)) and that commutators \([\bbF_{\mu\nu},\bbF_{\rho\sigma}]\) are small compared to \(\bbF\) in norm.
        Large commutators indicate intrinsically non-Abelian transport.
  \item \textbf{Trivial-connection sanity check.} In spatially uniform regions with a constant eigenframe,
        \(a_\mu\) should be (approximately) a pure gauge and the curvature should vanish up to numerical noise.
\end{enumerate}


\section{Symbol Dictionary}
\label{app:dictionary}

\begin{tabular}{ll}
\(\Omega\) & intrinsic domain (periodic cell or bounded cube) \\
\(\X(\x,t)\) & embedding map \(\Omega\to\R^4\) \\
\(\X_0(\x)\) & rest configuration \\
\(\uvec=\X-\X_0\) & displacement field \\
\(\vvec=\partial_t \uvec\) & velocity field \\
\(\omega_0\) & carrier angular frequency (Compton-based in project) \\
\(\mathbf{a}\) & complex amplitude \(\sqrt{\omega_0}\uvec + \ii \vvec/\sqrt{\omega_0}\) \\
\(\ket{\psi}\) & normalized ray \(\mathbf{a}/\|\mathbf{a}\|\) \\
\(a_\mu\) & U(1) Berry connection \(\ii\braket{u}{\partial_\mu u}\) \\
\(f_{\mu\nu}\) & U(1) curvature \(\partial_\mu a_\nu-\partial_\nu a_\mu\) \\
\(\bbA_\mu\) & Wilczek--Zee connection (matrix) \\
\(\bbF_{\mu\nu}\) & Wilczek--Zee curvature \(\partial_\mu\bbA_\nu-\partial_\nu\bbA_\mu-\ii[\bbA_\mu,\bbA_\nu]\) \\
\(Y_{\ell m}\) & spherical harmonic on \(S^2\) \\
\(\ell,m\) & harmonic degree/order (\(|m|\le \ell\)) \\
\(j_\ell\) & spherical Bessel function (first kind) \\
\(\alpha_{\ell n}\) & \(n\)-th positive root of \(j_\ell(\alpha)=0\) \\
\(a\) & chosen electron radius (intrinsic cavity radius) \\
\(k_{\ell n}\) & radial wavenumber \(\alpha_{\ell n}/a\) \\
\(\omega\) & electron seed angular frequency \(c_{\mathrm{wave}}k_{\ell n}\) \\
\(\x_c\) & electron center in intrinsic coordinates \\
\(\mathbf{p}_1,\mathbf{p}_2\) & orthonormal polarization vectors in \(\R^4\) \\
\(A\) & displacement amplitude scale used for normalization \\
\(\chi\) & dimensionless amplitude ratio \(A/a\) used in scaling arguments \\
\end{tabular}

% Uncomment if you use BibTeX:
% \bibliographystyle{unsrt}
% \bibliography{refs}

\end{document}
